Fix one row and abbreviate \(t=t_{n,s}\), \(\pi=\pi_{n,t}\), and \(M_a=\widehat m_{n,t-1,a}(X_{n,t})\).  Conditional on the past, the current covariate, and the two current potential outcomes, logging integrity, conditional treatment independence, and consistency from Definition~\(\mathrm{def:saco\text{-}early\text{-}stabilized\text{-}aipw\text{-}class}\) give
\[
\begin{aligned}
E[\widehat\phi_{n,t}\mid\mathcal B_{n,t},Y_{n,t}(0),Y_{n,t}(1)]
&=(Y_{n,t}(1)-M_1)-(Y_{n,t}(0)-M_0)+M_1-M_0\\
&=Y_{n,t}(1)-Y_{n,t}(0).
\end{aligned}
\tag{1}
\]
The current full-data triple is independent of the past and identically distributed within the row.  Taking conditional expectation of (1) given \(\mathcal F_{n,t-1}\) therefore yields
\[
E[\widehat\phi_{n,t}-\theta_{0,n}\mid\mathcal F_{n,t-1}]=0.
\tag{2}
\]
By the scored-order filtration in the definition, \(\mathcal G_{n,s-1}=\mathcal F_{n,t-1}\).  Hence (2) proves
\[
E[X_{n,s}\mid\mathcal G_{n,s-1}]=0.
\tag{3}
\]
This establishes the martingale-difference property in scored-set order; it does not treat the scored observations as independent.

We next derive, rather than assume, the fourth-moment envelope needed for the martingale argument.  Uniform overlap implies the pathwise bound
\[
|\widehat\phi_{n,t}|\le(1+\epsilon^{-1})
\{|Y_{n,t}(0)|+|Y_{n,t}(1)|+|M_0|+|M_1|\}.
\tag{4}
\]
For nonnegative \(u_1,\ldots,u_r\), convexity gives \((\sum_{j=1}^ru_j)^4\le r^3\sum_{j=1}^ru_j^4\).  Hence (4) and the two uniform source envelopes give
\[
E|\widehat\phi_{n,t}|^4
\le64(1+\epsilon^{-1})^4(2C_Y+2C_M).
\tag{5}
\]
Jensen's inequality and \(|x-y|^4\le8(|x|^4+|y|^4)\) also give
\[
|\theta_{0,n}|^4
\le E|Y_{n,t}(1)-Y_{n,t}(0)|^4\le16C_Y.
\tag{6}
\]
Combining (5)--(6),
\[
E|X_{n,s}|^4
\le C_X:=512(1+\epsilon^{-1})^4(2C_Y+2C_M)+128C_Y<\infty.
\tag{7}
\]
Every constant is common to the triangular-array class.  This is the required transfer from Saco's outcome and nuisance assumptions to the AIPW score; no bounded-increment premise has been introduced.

Apply Lemma~\(\mathrm{lem:early\text{-}measurable\text{-}l4\text{-}self\text{-}normalized\text{-}martingale\text{-}clt}\) with \(h_n=k_n\), \(\lambda_n=\Lambda_n\), and \(N_n=n_{\mathrm{eff}}\).  Equation (3) supplies the martingale premise, (7) supplies its fourth-moment premise, and the class definition supplies \(k_n/N_n\to0\), \(\Lambda_n\ge\lambda_->0\), and
\[
\frac{A_{n,+}^2}{N_n\Lambda_n}\longrightarrow1
\tag{8}
\]
uniformly.  The lemma therefore proves both
\[
\frac{Q_n}{N_n\Lambda_n}\longrightarrow1
\tag{9}
\]
in probability and the asserted uniform standard-normal limit for \(S_n/\sqrt{Q_n}\).  In particular \(P(Q_n>0)\to1\) uniformly.  Notice that (8) and \(\Lambda_n\ge\lambda_-\) also imply the retained source variance-growth condition with any fixed \(v_-<\lambda_-\) for all sufficiently large \(n\), but that numbered source condition was retained explicitly in the class rather than silently discarded.

It remains to pass from realized QV to Saco's feasible sample variance.  The definitions give the exact identities
\[
\widehat\theta_n-\theta_{0,n}=\frac{S_n}{N_n},\qquad
(N_n-1)\widehat V_n=Q_n-\frac{S_n^2}{N_n}.
\tag{10}
\]
Cauchy--Schwarz ensures \(0\le S_n^2\le N_nQ_n\).  The score-pivot limit implies \(S_n/\sqrt{Q_n}=O_P(1)\) uniformly, so on \(\{Q_n>0\}\),
\[
\frac{N_n\widehat V_n}{Q_n}
=\frac{N_n}{N_n-1}\left\{1-\frac1{N_n}\left(\frac{S_n}{\sqrt{Q_n}}\right)^2\right\}\longrightarrow1
\tag{11}
\]
in probability, uniformly.  Equations (9) and (11) imply \(P(\widehat V_n>0)\to1\).  On that event, (10)--(11) give
\[
Z_n^{\mathrm{feas}}
=\frac{S_n/\sqrt{Q_n}}{\sqrt{N_n\widehat V_n/Q_n}}\Rightarrow N(0,1)
\tag{12}
\]
uniformly.  Continuity of the standard normal distribution function upgrades (12) to the displayed uniform Kolmogorov conclusion.

Finally, with \(z_\alpha=\Phi^{-1}(1-\alpha/2)\), the event \(\{\theta_{0,n}\in\mathcal I_n(\alpha)\}\) agrees on \(\{\widehat V_n>0\}\) with \(\{|Z_n^{\mathrm{feas}}|\le z_\alpha\}\).  The exceptional event has probability tending uniformly to zero, and the normal limit has
\[
\Phi(z_\alpha)-\Phi(-z_\alpha)=1-\alpha.
\tag{13}
\]
The uniform distributional convergence in (12) and (13) prove the asserted Wald coverage.  Saco's ten numbered source conditions, Assumptions~3.5, 3.6, 3.2, 3.7, 3.8, 4.5, 4.8, 3.9, 3.10, and 5.11, are recorded explicitly in the preceding class definition.  The extra premise used above is precisely (8), so the theorem's class is their intersection with early QV stabilization, as claimed.